\section{SPECTRA-Bench}
\label{sec:bench}

\subsection{Interface, Devices and Testbenches}
Each task fixes a port list, a testbench and quantitative specifications. The designer (human or agent) returns one
subcircuit in a canonical netlist format restricted to SKY130 1.8\,V nFET/pFET devices~\cite{sky130pdk}, the
substrate PNP, resistors and capacitors; no ideal sources are allowed inside the design. Resistors are compiled to a
first-order SKY130 high-resistance poly model (PDK temperature coefficients), so temperature behaviour of
references, regulators and filters is not idealised. All circuits are simulated with ngspice-42~\cite{ngspice42}
using the PDK's BSIM4 models. Table~\ref{tab:families} lists the twelve families and their testbenches; each
testbench is a single batch run (or two to three runs) that extracts all metrics, and the op-amp testbench, for
example, measures open-loop gain/UGBW/PM through a DC-servo loop, CMRR, PSRR and input noise from the same run,
slew and settling in a unity-gain step, and the loaded linear output range in an inverting $\times10$
configuration. Every testbench was validated against hand analysis (e.g., an integrator's unity-gain frequency of
157.7\,kHz versus 159\,kHz expected).

\begin{table}[t]
\centering
\caption{Circuit families and measured metrics (SKY130, VDD\,=\,1.8\,V).}
\label{tab:families}
\scriptsize
\setlength{\tabcolsep}{3pt}
\begin{tabular}{@{}p{17mm}p{63mm}@{}}
\toprule
Family & Metrics extracted by the testbench \\
\midrule
Op-amp/OTA & gain, UGBW, PM, CMRR, PSRR, loaded swing, slew, settling, noise, power \\
Amplifier & gain, $f_{-3\mathrm{dB}}$, GBW, output DC, swing, power \\
TIA & transimpedance, $f_{-3\mathrm{dB}}$, input noise current, output DC, power \\
Current mirror & output-current error, $R_{out}$, compliance voltage \\
LDO & accuracy, load/line regulation, PSRR, $I_Q$, load-step under/overshoot, ringing \\
Reference & $V_{ref}$, TC ($-40$..125\,$^\circ$C), line reg., PSRR, start-up, power \\
CT comparator & delay at 10\,mV overdrive, offset, polarity, output levels, power \\
Clocked comp. & decision correctness, clock-to-output delay, energy/decision \\
VCO & frequency at 3 control voltages, tuning ratio, $K_{VCO}$, swing, power \\
Filter & passband gain, $f_c$ or $f_0$/$Q$, peaking, stop-band attenuation, power \\
Schmitt trigger & thresholds, hysteresis, output levels, power \\
Op-amp appl. & signed gains, closed-loop bandwidth, offset, integrator $f_u$ \\
\bottomrule
\end{tabular}
\end{table}


\subsection{Task Tiers and Feasibility Certificates}
Specifications are worthless if they are unattainable or trivially loose. For the \emph{certified tier} we wrote
\NTemplates{} expert parametric templates spanning the families, explored each with CMA-ES under random
scalarisations of the family's metrics (\NExploreEvals{} simulated designs in total), and selected reference
designs near the Pareto front. Specifications are the reference's metrics relaxed by a tier margin
(T1/T2/T3: e.g., $-10/-5/-2$\,dB for gains, $0.5/0.75/0.9\times$ for bandwidths, $2/1.33/1.1\times$ for power),
rounded conservatively, with 4, 6 or all profile metrics, respectively. Every task is re-simulated with its
reference before acceptance, so each certified task carries a \emph{feasibility certificate}. Robust tasks derive
targets from the reference's worst case over the full PVT grid and must hold on all 45 conditions.
Because certified specs come from our own templates, we add a \emph{designer tier} of \NDesigner{} tasks written
independently from textbook and practitioner targets~\cite{allen2012cmos,razavi2017design} (e.g., a 1.2\,V/10\,mA
on-chip LDO, a 1\,GHz dynamic comparator, a 1\,MHz Butterworth low-pass), whose feasibility is unknown a priori;
and \emph{oracle-unmet controls} that tighten certified tasks beyond any design in the exploration archives.
In total SPECTRA-Bench contains \NTasksTotal{} tasks (\NCertified{} certified, \NRobust{} of them robust).

\subsection{Metrics and Budgets}
A run succeeds only if an \emph{independent} re-simulation of the returned netlist meets every specification at TT
(nominal tasks) or at all 45 PVT conditions (robust tasks); final robust designs are additionally analysed with a
100-sample per-instance mismatch Monte Carlo using the PDK's Pelgrom coefficients~\cite{pelgrom1989matching}. All
methods share the same budget per run: at most 12 LLM calls and 300 (nominal) or 2{,}000 (robust)
\emph{condition-evaluations}, where evaluating one design at one PVT condition counts once regardless of how many
analyses the testbench needs; simulations spent on sensitivities or failed refinements are charged. We report
success with 95\% cluster-bootstrap confidence intervals over tasks, simulations and LLM calls to first success,
tokens, cost and wall-clock time.
